\section{One contour error comparison, with coordinate costs}
\label{sec:errors}

The bounded follow-up fixes one source-error model: a uniform
pointwise error in $B$ on one cleared circle. It compares two
explicit triangle/Cauchy allowances using the same data. It does
not compare exact feasible sets, source-acquisition costs or
canonical minimum-gap widths. Direct contour evaluation is
allowed to use the same combined kernel.

\begin{theorem}[Quadratic cancellation factor]
\label{thm:quadratic}
Use the exact block coordinates of Theorem~\ref{thm:general-flux}.
Let $\Gamma$ be the circle $|z-c|=A$, where
$|y_i|\le\rho<A$. Suppose $\widehat B$ satisfies
$\sup_\Gamma|B-\widehat B|\le\epsilon$. For the linear functional
$\E_J=2\av{BK_J}$, separate evaluation of the exterior values
and $b$ gives the valid allowance
\begin{equation}
 B_{\rm sep}=2A\epsilon\sum_i\frac{|J_i|}{A-|y_i|}
                      +\frac{2m|J|\epsilon}{A}.
 \label{eq:Bsep}
\end{equation}
The combined kernel gives the valid allowance
\begin{equation}
 B_{\rm comb}=\frac{2\epsilon}{A}
                  \sum_i\frac{|J_i||y_i|^2}{A-|y_i|},
 \qquad
 B_{\rm comb}\le\left(\frac\rho A\right)^2 B_{\rm sep}.
 \label{eq:Bcomb}
\end{equation}
The inequality is meaningful without dividing by $B_{\rm sep}$,
including cases where it is zero.
\end{theorem}
\begin{proof}
The contour length is $2\pi A$. From
\eqref{eq:exterior-Cauchy}, the error in each $R_i$ is bounded by
$A\epsilon/(A-|y_i|)$ and the error in $b$ by $\epsilon/A$.
Applying the triangle inequality to
$2\sum_iJ_iR_i-2mJb$ proves \eqref{eq:Bsep}.
The second form of \eqref{eq:Kgeneral} gives
\[
 \sup_\Gamma|K_J|\le\frac1{A^2}
                   \sum_i\frac{|J_i||y_i|^2}{A-|y_i|}.
\]
Multiply by $2A\epsilon$ to obtain $B_{\rm comb}$. Replacing
$|y_i|^2$ by $\rho^2$ proves
\[
 B_{\rm comb}\le(\rho/A)^2
       \left(2A\epsilon\sum_i\frac{|J_i|}{A-|y_i|}\right)
       \le(\rho/A)^2 B_{\rm sep}.
\]
\end{proof}

This is a strict factor less than one whenever the separate
allowance is positive and $\rho<A$. It is a comparison with the
specific allowance \eqref{eq:Bsep}; the latter is not asserted
to be the sharp independent-moment baseline. No superiority over
a direct method using \eqref{eq:Kgeneral} follows. This distinction
matches the parent's common-data comparison discipline
\cite{BakerCorrelation2026}.

\subsection{A quartet refinement and its asymmetry term}
\begin{proposition}[The next cancellation for a symmetric quartet]
\label{prop:cubic}
For $N=4$ and $J=C_4$, let $q_2=\sum_iJ_i y_i^2$. Then
\begin{equation}
 q_2=-\frac16p_2p_3,\qquad
 K_{C_4}(z)=\frac{q_2}{Z^3}
          +\frac1{Z^3}\sum_i\frac{J_i y_i^3}{Z-y_i}.
 \label{eq:cubic-kernel}
\end{equation}
Under Theorem~\ref{thm:quadratic}'s assumptions another valid
source-error allowance is
\begin{equation}
 B_{\rm cubic+asym}=\frac{2\epsilon|q_2|}{A^2}
       +\frac{2\epsilon}{A^2}
             \sum_i\frac{|J_i||y_i|^3}{A-|y_i|}.
 \label{eq:Bcubic}
\end{equation}
If the centered quartet is reflection-symmetric, $p_3=q_2=0$,
and
\begin{equation}
 B_{\rm cubic}\le(\rho/A)^3 B_{\rm sep}.
 \label{eq:cubic-factor}
\end{equation}
No parity condition is imposed on the source error $B-\widehat B$.
\end{proposition}
\begin{proof}
Since $J_i=-y_i^3+\beta p_2y_i+p_3/N$,
$q_2=-p_5+(\beta+1/N)p_2p_3$. For a centered quartet,
Newton identities give
$p_5=-b_2p_3-b_3p_2=5p_2p_3/6$, while
$\beta+1/N=5/12+1/4=2/3$. This proves the first identity.
Expanding one more term in \eqref{eq:Kgeneral} proves the second.
The circle-length bound gives \eqref{eq:Bcubic}. When $p_3=0$,
replace $|y_i|^3$ by $\rho^3$ and compare with the first term of
\eqref{eq:Bsep}, as in Theorem~\ref{thm:quadratic}.
\end{proof}

Approximate symmetry alone cannot set $q_2$ to zero. The first
term of \eqref{eq:Bcubic} is its explicit payment; even that
allowance need not improve \eqref{eq:Bcomb} in every asymmetric
configuration.

\subsection{Uncertain block coordinates}
The preceding factors keep the coordinates exact. A root enclosure
must pay for the kernel's coordinate dependence as well. The
following conservative bound supplies one such payment for $C_4$.

\begin{proposition}[A coordinate allowance for the fourth-order kernel]
\label{prop:coordinates}
Let $x$ and $\widehat x$ be two labeled $N$-tuples, $N\ge4$, with
$\max_i|x_i-\widehat x_i|\le e$. On the straight homotopy between
them, assume every centered root has modulus at most $\rho$.
Let one fixed circle $\Gamma$ have radius $A_0$. Suppose for every
homotopy parameter and every $z\in\Gamma$ that
\begin{equation}
 |z-x_i|\ge\delta>0,\qquad |z-c|\ge\delta>0.
 \label{eq:coordinate-clearance}
\end{equation}
For $K=K_{C_4}$,
\begin{equation}
 \sup_\Gamma|K(x)-K(\widehat x)|\le G,
 \quad
 G=Ne\left(\frac{20\rho^2}{\delta}
          +\frac{20\rho^3}{\delta^2}
          +\frac{6\rho^4}{\delta^3}\right).
 \label{eq:G}
\end{equation}
If $|\widehat B|\le B_*$ and $|B-\widehat B|\le\epsilon$ on
that circle, the corresponding exact contour functionals obey
\begin{equation}
 |\E_4(B,x)-\E_4(\widehat B,\widehat x)|
 \le 2A_0\left[\epsilon\sup_\Gamma|K(\widehat x)|
                         +(B_*+\epsilon)G\right].
 \label{eq:joint-coordinate}
\end{equation}
Here $\E_4(B,x)$ denotes the functional in
\eqref{eq:contour-general}; its identification with the physical
flux uses the actual counted block $x$.
\end{proposition}
\begin{proof}
Differentiating \eqref{eq:psi} in a coordinate gives
\begin{equation}
 \partial_{x_j}\psi_i
  =(3y_i^2-\beta p_2)(\delta_{ij}-1/N)
       -2\beta y_i y_j-\frac3N(y_j^2-p_2/N).
 \label{eq:psi-jacobian}
\end{equation}
Since $\beta N<2$, the sums of absolute values of these three
terms over $j$ are at most $10\rho^2$, $4\rho^2$ and
$6\rho^2$, respectively. Thus
$\sum_j|\partial_{x_j}\psi_i|\le20\rho^2$.
Also $|J_i|=|\psi_i|\le4\rho^3$ and, from \eqref{eq:C4},
$|C_4|\le3N\rho^4/4$. Along the homotopy, these imply
\[
 |\dot J_i|\le20\rho^2 e,\qquad
 |\dot C_4|\le4N\rho^3 e,\qquad |\dot c|\le e.
\]
Differentiate
$K=\sum_iJ_i/(z-x_i)-4C_4/(z-c)^2$ along the homotopy:
\[
 \dot K=\sum_i\left(\frac{\dot J_i}{z-x_i}
                       +\frac{J_i\dot x_i}{(z-x_i)^2}\right)
       -\frac{4\dot C_4}{(z-c)^2}
       -\frac{8C_4\dot c}{(z-c)^3}.
\]
The four contributions are bounded by
$20N\rho^2e/\delta$, $4N\rho^3e/\delta^2$,
$16N\rho^3e/\delta^2$ and $6N\rho^4e/\delta^3$.
Integrate over the unit homotopy interval to obtain \eqref{eq:G}.
Finally use the exact splitting
\[
 BK(x)-\widehat B K(\widehat x)
  =(B-\widehat B)K(\widehat x)+B\bigl(K(x)-K(\widehat x)\bigr)
\]
and $|B|\le B_*+\epsilon$ in the circle-length bound. This also
pays the mixed field/coordinate error.
\end{proof}

If the fixed circle is centered at the nominal mean, the cancellation
bounds apply to $K(\widehat x)$. Its distance to the actual mean
and roots must still satisfy \eqref{eq:coordinate-clearance} along
the whole homotopy. The term involving $G$ can erase the nominal
source-error gain. A numerical contour quadrature introduces a
separate error for the integrand actually integrated; arithmetic
enclosures and any integration of \eqref{eq:integrated-C4} in time
must be paid too. Counts, clearance, and between-time control are
additional premises, not consequences of the radius factor.

\subsection{Two fixed exact arithmetic checks}
The bounded check uses two six-root polynomial initial slices,
and no canonical source acquisition. For the block
$(-2,-1,1,2)$ with exterior $(-5,5)$, exact rational evaluation gives
\begin{equation}
 C_4=\frac{23}{12},\quad C_4'=-\frac{55}{42},\quad
 b=-\frac2{25},\quad \E_4=-\frac{29}{350}.
 \label{eq:symmetric-example}
\end{equation}
The same block without the exterior has $C_4'=0$. At radius
$A=4$ and unit source-error scale, the specified allowances are
\begin{equation}
 B_{\rm sep}=\frac{421}{18},\quad
 B_{\rm comb}=\frac{31}{18},\quad
 B_{\rm cubic}=\frac{43}{72}.
 \label{eq:example-allowances}
\end{equation}
Thus their exact ratios to the separate allowance are $31/421$
and $43/1684$. These are allowance ratios on exact finite geometry,
not measured errors or canonical width reductions.

For the asymmetric block $(-3,-1,2,5)$ with exterior $(-8,9)$,
the center is $3/4$. At radius $A=6$ the roots have
$\rho=17/4$, and exact arithmetic gives
\begin{align}
 p_2&=147/4,&p_3&=165/8,&C_4&=861/128,\notag\\
 q_2&=-8085/64,&C_4'&=-6847/7280,&
 \E_4&=7306487/13213200.
 \label{eq:asymmetric-example}
\end{align}
This checks the nonzero asymmetry payment. The execution record
\cite{BakerAffineChecks2026} contains both sets of allowance values
and a quartic collision at $T=1/20$, $b=2$. In each six-root case
it computes the derivative from all root velocities, the exterior
balance, and the exterior residues of the combined contour kernel
independently in exact rational arithmetic. These checks confirm
the finite algebra; the general theorems rest on their proofs.
